\IEEEraisesectionheading{\section{Introduction}\label{sec:introduction}}

\IEEEPARstart{S}{tress} can have detrimental effects on the human body, ranging from depression, anxiety and social withdrawal to sleep deprivation and obesity. Stress has a direct impact on the cardiovascular system \cite{mcewen2017neurobiology}. A wide variety of methods exist to monitor these cardiovascular signals and to track stress levels. The measurement of pulse waves is called plethysmography. Photoplethysmography (PPG) utilizes light reflectance to detect physiological changes. A traditional method for PPG is contact-based sensors, such as pulse oximeters \cite{verkruysse2008remote}. Remote Photoplethysmography (rPPG) uses standard video cameras to monitor the skin color and filter for subtle skin-color variations that are caused by flowing blood. Based on this information, the heart rate can be calculated \cite{verkruysse2008remote}. These measurements can be directly reported back to the user using biofeedback. Biofeedback describes a therapeutic concept in which a user receives real-time information about their physiological processes, such as respiration or heart rate. This data is typically made accessible through an output device, such as a visual display. Alternatively, the information can be translated into different forms, such as auditory signals. The primary goal of biofeedback is to make these biological processes perceptible, enabling the user to gain conscious awareness and control over functions that were previously automatic \cite{giggins2013biofeedback}.


% hier nocheinma über die übersicht überarbeiten 
The objective of this paper is to evaluate such a camera-based stress tracking system combined with musical biofeedback. The goal is to analyze if such a system can be used to lower stress levels and regulate consciousness. In the second chapter theoretical foundations such as stress, consciousness and rPPG techniques are investigated. In the third chapter, case studies will be analyzed. The fourth chapter evaluates whether such a system is theoretically feasible.